\section{The tree circuit as a multiscale network}
\label{sec:ldp_tree_mera}

The paragraph ``Connection to MERA'' of~\cite{Malz2024Preparation} reads the
tree circuit of Theorem~\ref{thm:ldp_tree_factorization} as a finite-range
multiscale entanglement renormalization ansatz with $O(\log\log N)$ layers: the
isometries $V^{(j)}$ are the isometries of the network, and all disentanglers
are the identity except those of the first layer, which prepare the
fixed-point state. This section states the network, identifies the
approximating state with its state, and counts the layers. Throughout, the
two-site blocked tensor of $A$ is injective and the chain length is
$N=M2^{k+1}$; both restrictions are recorded
in~\cite{gap:mswc24_tree_mera_scope}.

\begin{definition}[Binary finite-range network]
    \label{def:ldp_binary_mera}
    \lean{MPSTensor.binaryTreeMatrix, MPSPreparation.BinaryMERA,
        MPSPreparation.BinaryMERA.treeMatrix,
        MPSPreparation.BinaryMERA.topPair,
        MPSPreparation.BinaryMERA.amplitude,
        MPSPreparation.BinaryMERA.state,
        MPSPreparation.BinaryMERA.state_apply}
    \leanok
    \uses{def:ldp_approximating_state, thm:ldp_tree_factorization}
    Let $n,D\ge1$ and $k\ge0$. A \emph{binary finite-range network} with $k+1$
    isometry layers consists of a unitary $u$ on $\C^D\otimes\C^D$, isometries
    $W_0,\dots,W_{k-1}\colon\C^{D^2}\to\C^{D^2}\otimes\C^{D^2}$ and an isometry
    $W\colon\C^{D^2}\to\C^n\otimes\C^n$. Its tree map is
    \begin{align}
        \mathcal W
         & =W^{\otimes2^k}\,W_0^{\otimes2^{k-1}}\cdots W_{k-1}
        \colon\C^{D^2}\to(\C^n)^{\otimes2^{k+1}},
        \label{eq:ldp_binary_mera_tree}
    \end{align}
    where the Kronecker powers act on the regrouping of the $2^{k+1}$ sites
    into neighbouring pairs, then pairs of pairs, and so on, as in
    (\ref{eq:ldp_tree_factorization}). On $M$ top sites
    $\C^{D^2}=L_i\otimes R_i$, the state of the network on $N=M2^{k+1}$ sites
    is
    \begin{align}
        \mathcal W^{\otimes M}\bigotimes_{i=1}^{M}
        u_{R_iL_{i+1}}\ket{0}_{R_i}\ket{0}_{L_{i+1}},
        \label{eq:ldp_binary_mera_state}
    \end{align}
    with $L_{M+1}=L_1$. Every isometry maps one site to two neighbouring
    sites and every unitary acts on two neighbouring legs, so the network has
    range two and bond dimension $D^2$ in every layer; the disentanglers below
    the first layer are the identity.
\end{definition}

\begin{lemma}[The network state is a unit vector]
    \label{lem:ldp_binary_mera_norm}
    \lean{MPSTensor.isIsometry_binaryTreeMatrix,
        MPSPreparation.BinaryMERA.isIsometry_treeMatrix,
        MPSPreparation.BinaryMERA.sum_star_topPair_mul,
        MPSPreparation.BinaryMERA.norm_state}
    \leanok
    \uses{def:ldp_binary_mera, thm:prep_pair_product_norm}
    The tree map (\ref{eq:ldp_binary_mera_tree}) of a binary finite-range
    network is an isometry, and the state (\ref{eq:ldp_binary_mera_state})
    is a unit vector.
\end{lemma}
\begin{proof}
    \leanok
    A Kronecker power of an isometry is an isometry, and so is a product of
    isometries, so induction on $k$ shows that $\mathcal W$ is an isometry.
    The pair $u\ket0\ket0$ is a unit vector because $u$ is unitary, so the
    product of pairs is a unit vector
    (Theorem~\ref{thm:prep_pair_product_norm}), and $\mathcal W^{\otimes M}$
    preserves its norm.
\end{proof}

\begin{lemma}[The layers of the tree circuit are isometries]
    \label{lem:ldp_tree_layers_isometries}
    \lean{MPSTensor.pairPosTensor, MPSTensor.treeLayers,
        MPSTensor.binaryTreeMatrix_treeLayers,
        MPSTensor.isInjective_of_isInjective_rotatePhysical,
        MPSTensor.isInjective_blockTensor_two,
        MPSTensor.isInjective_blockTensor_pairPosTensor,
        MPSTensor.isInjective_blockTensor_iterate_pairPosTensor,
        MPSTensor.isIsometry_treeLayers}
    \leanok
    \uses{thm:ldp_tree_factorization, thm:ldpolar_tensor_injective,
        thm:block_tensor_rotate_physical}
    Let the two-site blocked tensor of $A$ be injective, and let $V^{(j)}$ and
    $T_j$ be as in Theorem~\ref{thm:ldp_tree_factorization}. Then for every
    $j\ge1$ the two-site blocked tensor of $T_{j-1}$ is injective, so each
    $V^{(j)}$ is an isometry, and the tree map (\ref{eq:ldp_binary_mera_tree})
    with $W=V^{(1)}$ and $W_i=V^{(i+2)}$ is the product of the layers in
    (\ref{eq:ldp_tree_factorization}) for $q=2^{k+1}$.
\end{lemma}
\begin{proof}
    \leanok
    If a tensor $Y$ rotated by a matrix on its physical leg is injective, so
    is $Y$, because the matrices of the rotated tensor are combinations of those
    of $Y$. The two-site blocked tensor of an injective tensor is injective.
    Blocking $\mathcal B^{(1)}=V^{(1)}P_1$ twice gives $(V^{(1)})^{\otimes2}$
    applied to $T_1$ blocked twice
    (Theorem~\ref{thm:block_tensor_rotate_physical}), so the two-site blocked
    tensor of $T_1$ is injective, and induction on $j$ gives the claim for
    every $T_{j-1}$. Its isometric factor is an isometry by
    Theorem~\ref{thm:ldpolar_tensor_injective}. The identification of the
    tree map with the product of the layers is the recursion defining both.
\end{proof}

\begin{definition}[The tree network of a tensor]
    \label{def:ldp_tree_mera}
    \lean{MPSPreparation.treeMERA,
        MPSPreparation.exists_disentangler_fixedPointPair}
    \leanok
    \uses{def:ldp_binary_mera, lem:ldp_tree_layers_isometries,
        def:prep_pair_product_state, thm:prep_pair_norm}
    Let the two-site blocked tensor of $A$ be injective and let $u$ be a unitary
    on $\C^D\otimes\C^D$. The \emph{tree network} of $A$ with $k+1$ isometry
    layers is the binary finite-range network with isometries
    $W=V^{(1)}$, $W_i=V^{(i+2)}$ for $0\le i<k$, and unitary $u$. For
    $\rho\ge0$ with $\tr\rho=1$ there is a unitary $u$ with
    $u\ket0\ket0=\ket\omega$, the pair of the fixed-point state.
\end{definition}

\begin{theorem}[The approximating state is a network state]
    \label{thm:ldp_approximating_state_tree_mera}
    \lean{MPSPreparation.approximatingMPVState_eq_state_treeMERA}
    \leanok
    \uses{def:ldp_tree_mera, def:ldp_approximating_state,
        thm:ldp_tree_factorization, lem:ldp_approximating_state_unit_norm}
    Let the two-site blocked tensor of $A$ be injective, let $\rho\ge0$ with
    $\tr\rho=1$, and let $u$ be a unitary with $u\ket0\ket0=\ket\omega$. For
    $q=2^{k+1}$ and $M\ge1$, the approximating state
    $\ket{\widetilde\phi_N}=V^{\otimes M}\ket\Omega$ on $N=Mq$ sites
    (Definition~\ref{def:ldp_approximating_state}) is the state of the tree
    network of $A$ with $k+1$ isometry layers and unitary $u$.
\end{theorem}
\begin{proof}
    \leanok
    The blocked tensor over $q=2^{k+1}$ sites is injective, so the
    approximating state is the unnormalized one
    (Lemma~\ref{lem:ldp_approximating_state_unit_norm}). Its isometry $V$ is
    the product of the layers by Theorem~\ref{thm:ldp_tree_factorization},
    hence the tree map by Lemma~\ref{lem:ldp_tree_layers_isometries}, and the
    pairs $u\ket0\ket0$ are the pairs $\ket\omega$ of $\ket\Omega$.
\end{proof}

\begin{theorem}[Normal states are networks with $O(\log\log(N/\varepsilon))$ layers]
    \label{thm:ldp_tree_mera_error}
    \lean{MPSPreparation.exists_state_treeMERA_approximationError_le,
        MPSPreparation.exists_state_treeMERA_approximationError_le_of_lt}
    \leanok
    \uses{thm:ldp_approximating_state_tree_mera,
        thm:ldp_approximation_error_normal, def:correlation_length}
    Let $A$ be normal in the gauge~(\ref{eq:ldp_normal_gauge}), let $0<t<1$
    bound the moduli of the eigenvalues of $\E_A$ other than $1$, with
    $\xi=-1/\log t$, and let the two-site blocked tensor of $A$ be injective.
    There is $C>0$ such that for every unitary $u$ with
    $u\ket0\ket0=\ket\omega$, every $\varepsilon>0$, and every $k\ge0$ and
    $M\ge1$ with $N=M2^{k+1}$ and
    \begin{align}
        x & =\xi\log(CN/\varepsilon)\le2^{k+1},
        \label{eq:ldp_tree_mera_threshold}
    \end{align}
    the state $\ket\psi$ of the tree network of $A$ with $k+1$ isometry layers
    satisfies $\varepsilon(\psi,\phi_N)\le\varepsilon$. If moreover
    $2^k<x$, that is, $k$ is the least exponent with
    (\ref{eq:ldp_tree_mera_threshold}), then $k<\log_2x$, so the network has
    fewer than $\log_2(\xi\log(CN/\varepsilon))+1$ isometry layers. This is
    the containment of normal translation-invariant matrix product states in
    finite-range networks with $O(\log\log N)$ layers, within the
    approximation error, of~\cite[paragraph ``Connection to
        MERA'']{Malz2024Preparation}, for $N=M2^{k+1}$ and tensors whose
    two-site blocked tensor is injective~\cite{gap:mswc24_tree_mera_scope}.
\end{theorem}
\begin{proof}
    \leanok
    By Theorem~\ref{thm:ldp_approximating_state_tree_mera} the network state
    is the approximating state, and Theorem~\ref{thm:ldp_approximation_error_normal}
    at $\gamma=1/2$ gives $C>0$ with
    $\varepsilon(\widetilde\phi_N,\phi_N)\le CMe^{-q/\xi}$. From
    (\ref{eq:ldp_tree_mera_threshold}) and $M\le N$,
    $\log(CM/\varepsilon)\le q/\xi$, so $CMe^{-q/\xi}\le\varepsilon$. The
    bound on $k$ is $2^k<x$ read through $\log_2$.
\end{proof}
